% TikZ figure: oldforms vs newforms
% Shows the inclusion/projection between levels N and M|N
\begin{figure}[ht]
\centering
\begin{tikzpicture}[
  box/.style={draw, rounded corners=4pt, minimum width=2.8cm, minimum height=1cm, font=\small, align=center},
  arrow/.style={->, thick, >=stealth}
]

% Level M box
\node[box, fill=green!15] (SM) at (0, 3) {$\Sk_k(M)$\\(低级别)};

% Map down to level N
\node[box, fill=blue!10] (SN) at (4, 3) {$\Sk_k(N)$\\(高级别 $M \mid N$)};

% Oldforms inside SN
\node[box, fill=orange!15, dashed, draw=orange] (old) at (4, 1) {旧形式 $\Sk_k^{\mathrm{old}}(N)$};

% Newforms inside SN
\node[box, fill=red!10, dashed, draw=red] (new) at (7.5, 3) {新形式 $\Sk_k^{\mathrm{new}}(N)$};

% Arrows
\draw[arrow] (SM) -- node[above, font=\small] {$f(\tau) \mapsto f(d\tau)$, $d = N/M$} (SN);
\draw[arrow, dashed] (SN.south) -- (old.north) node[midway, right, font=\small] {投影};

% Decomposition
\draw[arrow, double] (SN.east) -- node[above, font=\small] {正交分解} (new.west);

% Horizontal brace for decomposition
\node[font=\small] at (5.5, 4.0) {$\Sk_k(N) = \Sk_k^{\mathrm{old}}(N) \oplus \Sk_k^{\mathrm{new}}(N)$};

% Annotations
\node[font=\scriptsize, text=orange!80!black, align=center] at (4, 0.1)
  {由低级别 $M \mid N, M < N$ 的形式\\经 $f \mapsto f(d\tau)$ 生成};
\node[font=\scriptsize, text=red!80!black, align=center] at (7.5, 2)
  {与所有旧形式正交\\满足强重数一定理};

\end{tikzpicture}
\caption{$\Sk_k(N)$ 的旧形式—新形式正交分解。旧形式由低级别升上来，新形式是"真正属于级别 $N$"的部分。}
\label{fig:oldforms-newforms}
\end{figure}
